\documentclass[10pt,a4paper]{amsart}
\usepackage[margin=19mm]{geometry}
\usepackage{amsmath,amssymb,mathtools}
\usepackage[scr=rsfs,cal=euler]{mathalfa}
\usepackage{xfrac}
\usepackage[unicode,hidelinks,bookmarks=false]{hyperref}
\newcommand{\ie}{i.e.\ }
\newcommand{\cf}{cf.\ }
\newcommand{\resp}{respectively\ }
\newcommand{\Subsection}{\subsection}
\providecommand{\nobreakdash}{\nobreak\hbox{-}}
\setlength{\emergencystretch}{2em}
\multlinegap0pt
\usepackage{bm}
\usepackage{bbm}
\usepackage{dsfont}
\usepackage{xspace}
\usepackage{cancel}
\usepackage{mathtools}
\usepackage[shortlabels]{enumitem}
\usepackage{amsthm}
\makeatletter
\@ifundefined{thm}{\newtheorem{thm}{Theorem}}{}
\@ifundefined{lem}{\newtheorem{lem}[thm]{Lemma}}{}
\makeatother
\def\N{\mathbb{N}}
\def\R{\mathbb{R}}
\title[Intrinsic Ultracontractivity for a class of Schroedinger Sem]{Intrinsic Ultracontractivity for a class of Schroedinger Semigroups in $L^{2}(\mathbb{R}^{n})$ by Logarithmic Sobolev inequalities}
\author{Christoph Schwerdt, Alexander Mill, Dirk Hundertmark}
\date{Source: arXiv:2602.04685v1 (CC BY 4.0)}
\begin{document}
\maketitle

\noindent\emph{Source and licence.} Intrinsic Ultracontractivity for a class of Schroedinger Semigroups in $L^{2}(\mathbb{R}^{n})$ by Logarithmic Sobolev inequalities. Version arXiv:2602.04685v1 (CC BY 4.0), distributed under the Creative Commons Attribution 4.0 International licence, as recorded in the arXiv licence metadata for this exact version and shown on the abstract page https://arxiv.org/abs/2602.04685v1. Full rights notice: licenses/arxiv-2602.04685-CC-BY-4.0.txt.\par\smallskip

\noindent\textit{Editorial scope.} Counted unit: the Thm whose source label is \texttt{None} in the article (source0.tex lines 1140--1142, source label \texttt{None}), delivered together with the article's own complete original proof of it (source0.tex lines 1144--1222, one \texttt{proof} environment of 79 lines). No mathematics is written, repaired or re-derived: statement and proof are the article's own text line for line. Required context retained uncounted: CharacterizationIrreducibilty (lines 1128--1136). Every definition, display and label of those lines is retained unchanged. Omitted from this package and not redistributed: 1--1127, 1137--1139, 1143--1143, 1223--1311. Nothing inside a retained excerpt is omitted or altered: every retained body is delivered exactly as the article's lines are written, so no omission receipt of any kind accompanies this package. Citations: the retained excerpts carry no citation command at all, so no bibliography entry of the article is transcribed and no reference is redistributed. Reference adaptation: none was needed and none was made; every reference inside the delivered unit resolves inside it, so no unresolved reference or citation is produced. Printed slips kept literal: no printed word, symbol, hypothesis or number of the counted statement or of its proof was altered. Layout and class: the article's own class is \texttt{article} with packages including inputenc, fontenc, graphicx, amsmath, amsfonts, mathrsfs, amsthm, amssymb, color, enumerate, geometry, hyperref, bbm, breakurl; the wrapper is the lane's pinned \texttt{amsart} editorial wrapper at 10pt on A4, which restates the theorem-like environments the retained text needs and adds the article's own macro definitions that the retained text needs (\texttt{\textbackslash N}, \texttt{\textbackslash R}), with extra wrapper packages (bm, bbm, dsfont, xspace, cancel, mathtools, enumitem). The delivered numbering is the unit's own. Assets: no retained span contains a figure environment, a graphics-inclusion command, a PDF-page inclusion, a table or a TikZ picture; no image file is copied into this package, the package declares no asset dependency and no image file is redistributed with it. Licence: version arXiv:2602.04685v1 of this article is distributed under the Creative Commons Attribution 4.0 International licence, as recorded in the arXiv licence metadata for this exact version and shown on the abstract page https://arxiv.org/abs/2602.04685v1. A full rights notice is delivered at licenses/arxiv-2602.04685-CC-BY-4.0.txt. Characters: every retained excerpt is pure ASCII, so the delivered engine, XeLaTeX with its default Latin Modern text font, reports no missing character.
\begin{lem}\label{CharacterizationIrreducibilty}
If $h$ satisfies $h(u, v) = 0$ for functions $u, v \in D(h)$ having disjoint 
supports then the following statements are equivalent:
 \begin{enumerate}[i.)]
\item For every $t>0$ the operators $\mathrm{e}^{-tH}$ are positivity improving.
\item If $\Omega \subseteq \R^{n}$ satisfies $1_{\Omega} \ u \in D(h)$ for every 
$u \in D(h)$ then either $\Omega$ or $\R^{n} \setminus \Omega$ is a Lebesgue null set.\\
\end{enumerate} 
\end{lem} 

\begin{thm}
For every $t>0$ the operators $\mathrm{e}^{-tH}$ is positivity improving.\\
\end{thm}

\begin{proof}
We prove by contradiction and assume that $\{ \mathrm{e}^{-tH} \ | \  t \geq 0 \}$ is not positivity improving. 
By Lemma \ref{CharacterizationIrreducibilty} there is a set $\Omega \subseteq \R^{n}$ such that  
$1_{\Omega} u$  is contained in $D(h)$ for every $u \in D(h)$ but neither $\Omega$ nor 
$\Omega^{C}  = \R^{n} \setminus \Omega$ is a Lebesgue null set. We denote the Lebesgue measure in 
$\R^{n}$ by $\lambda$.\\

\begin{enumerate}[i.)]
\item We show that there exists a $x_{0} \in \R^{n}$ such that neither $\Omega \cap B(x_{0},r)$ nor 
$\Omega^{C} \cap B(x_{0},r)$ is a Lebesgue null set for any radius $r>0$.\\

We assume that this claim would be false. Then for every $x \in \R^{n}$ there is a radius $r=r(x) > 0$ such that either
$\Omega \cap B(x,r)$ or $\Omega^{C} \cap B(x,r)$ is a Lebesgue null set. We define the sets $\Omega_{1}$ by
$$
\Omega_{1} = \{ x \in \R^{n} \ | \ \exists r > 0 \ : \ \lambda(\Omega \cap B(x,r)) = 0 \}
$$
and $\Omega_{2}$ by $\{ x \in \R^{n} \ | \ \exists r > 0 \ : \ \lambda(\Omega^{C} \cap B(x,r)) = 0 \}$. Then 
$\R^{n} = \Omega_{1} \cup \Omega_{2}$ is true. Furthermore, mind that both sets 
are disjoint. If there would exists a $x \in \Omega_{1} \cap \Omega_{2}$ then
$$
0  = \lambda(\Omega \cap B(x,r)) + \lambda(\Omega^{C} \cap B(x,r)) = \lambda \bigl( B(x,r) \bigr) > 0
$$
would follow for $r>0$ chosen appropriately small enough. So, $\Omega_{1} \cap \Omega_{2}$ must be empty. 
It is also clear that $\Omega_{1}$ and $\Omega_{2}$ are 
both open sets in $\R^{n}$. But then either $\R^{n} = \Omega_{1}$ or $\R^{n} = \Omega_{2}$ must be true since $\R^{n}$
is connected.\\

Let $\R^{n} = \Omega_{1}$ be true. We choose a compact set $K \subseteq \Omega$ such that $\lambda(K) > 0$. This is possible
since $\Omega$ is not a pure point set in $\R^{n}$. For every $x \in K$ there exists a $r = r(x) > 0$ such that 
$$
\lambda (B(x,r) \cap \Omega)=0
$$ 
since $K \subseteq \Omega \subseteq \R^{n}=\Omega_{1}$. By the compactness of $K$ we conclude that 
$$
K \subseteq B(x_{1},r_{1}) \cup \dots \cup B(x_{m},r_{m}) 
$$
holds for some $x_{1}, \dots, x_{m} \in K$ where $r_{j}$ are given by $\Omega_{1}$. But then
$$
0 < \lambda(K) = \lambda(K \cap \Omega) \leq \sum_{j=1}^{m} \lambda \bigl( B(x_{j},r_{j}) \cap \Omega \bigr) = 0
$$
is obviously a contradiction. Using the same argumentation for $\R^{n} = \Omega_{2}$ leads to a contradiction as well. 
Hence there must exist a $x_{0} \in \R^{n}$ such that neither $\Omega \cap B(x_{0},r)$ nor $\Omega^{C} \cap B(x_{0},r)$ 
is a Lebesgue null set for any radius $r>0$.
\item Let $u \in C_{c}^{\infty}(\R^{n}) \subseteq D(h)$ be a function such that $u(x_{0}) = 1$. 
So $1_{\Omega} u$ is contained in $D(h)$ as we mentioned in the very beginning of this proof. 
In particular $1_{\Omega} u$ is an element of $H^{1}(\R^{n})$ with 
$\partial_{j} (1_{\Omega} u) = 1_{\Omega} \partial_{j}u$. But this last equation needs further explanation. First remember
that $\Omega$ is not a Lebesgue null set. Therefore we choose $v \in C_{c}^{\infty}(\Omega)$. So
\begin{align*}
& \int_{\Omega} (\partial_{j} u)(x) \overline{v(x)} \ \mathrm{d}x = - \int_{\Omega} u(x) \overline{(\partial_{j} v)(x)} \ \mathrm{d}x
 =  - \int_{\R^{n}} (1_{\Omega} u)(x) \overline{(\partial_{j} v)(x)} \ \mathrm{d}x \\
& = \int_{\R^{n}} \partial_{j}(1_{\Omega} u)(x) \overline{v(x)} \ \mathrm{d}x
 = \int_{\Omega} \partial_{j}(1_{\Omega} u)(x) \overline{v(x)} \ \mathrm{d}x
\end{align*}
is true where the last equation is due to the support of $v$ in $\Omega$. Hence we see that $\partial_{j} u$ is equal to 
$\partial_{j}(1_{\Omega} u)$ in $\Omega$ again due to the fundamental lemma in the calculus of variations. Now we remember 
that $\Omega^{C}$ is also not a Lebesgue null set. We
use the similar arguments for $v \in C_{c}^{\infty}(\Omega^{C})$. So
\begin{align*}
0 & = - \int_{\R^{n}} (1_{\Omega} u)(x) \overline{(\partial_{j} v)(x)} \ \mathrm{d}x = 
\int_{\R^{n}} \partial_{j}(1_{\Omega} u)(x) \overline{v(x)} \ \mathrm{d}x \\
& = \int_{\Omega^{C}} \partial_{j}(1_{\Omega} u)(x) \overline{v(x)} \ \mathrm{d}x 
\end{align*}
holds due to the support of $v$ in $\Omega^{C}$. Therefore $\partial_{j}(1_{\Omega} u)$ is $0$ in $\Omega^{C}$. So, 
in total we have shown that $\partial_{j} (1_{\Omega} u) = 1_{\Omega} \partial_{j}u$ is actually true.\\
Since $u \in C_{c}^{\infty}(\R^{n})$ we infer that $1_{\Omega} u$ is contained in the Sobolev spaces 
$W^{1,p}(\R^{n})$ for every $p \geq 2$. Sobolev embedding theorems state that $1_{\Omega} u$ has a 
continuous representative on $\R^{n}$. We treat $1_{\Omega} u$ as a continuous function without minding mathematical rigor.\\

Remember that for $x_{0}$ we have shown that neither $\Omega \cap B(x_{0},\frac{1}{k})$ nor $\Omega^{C} \cap B(x_{0},\frac{1}{k})$
are Lebesgue null sets for every $k \in \N$. Hence we choose two sequences by
$x_{k} \in \Omega \cap B(x_{0},\frac{1}{k})$ and $y_{k} \in \Omega^{C} \cap B(x_{0},\frac{1}{k})$. Then 
$$ 
|(1_{\Omega} u)(x_{k}) -  (1_{\Omega} u)(y_{k})| = |u(x_{k})|
$$
converges to $|u(x_{0})|$ for $k \to \infty$. But since $u(x_{0})$ is equal to $1$ by the choice of $u$ this contradicts the 
continuity of $1_{\Omega} u$. Hence our assumption leads to a contradiction which proves the claim.
\end{enumerate}
\end{proof}

\end{document}
